To check exactness of the sequence, we can work locally on @@M0@@ and use the sections defined in the statement of 2.5. We then have a basis of @@M1@@ formed by the images of the @@M2@@, a basis of @@M3@@ formed by the images of the @@M4@@, and a basis of @@M5@@ formed by the images of the @@M6@@. It is then clear that @@M7@@ is surjective, and that its kernel consists of the elements whose coordinates on the @@M8@@ and the @@M9@@ vanish, and whose coordinates on @@M11@@ and @@M12@@ are equal for every @@M10@@. Consequently the sequence is also exact in the middle. The three @@M13@@-Modules involved are locally free of finite type; since @@M14@@ has rank @@M15@@ by the proof of i), we see immediately that the sequence is exact.

\sgaStableTarget{sga6.e07.para.2-6}\subsection*{2.6.}
We now suppose that @@M16@@ is a scheme having an ample sheaf. Every closed subscheme @@M17@@ of @@M18@@ then has an ample sheaf. We know that in this case ``global resolutions exist'', and that the @@M19@@ formed respectively with locally free sheaves of finite type and with perfect complexes are identical (II, 2.2.9). If moreover the immersion @@M20@@ is regular, it is a perfect morphism, and we can define an additive homomorphism @@M21@@ (IV 2.12) by associating with the class of a perfect complex @@M22@@ the class of the complex @@M23@@ (@@M24@@ being an exact functor).

\textbf{Proposition 2.7.} \uline{Let} @@M25@@ \uline{be a scheme having an ample sheaf}, @@M26@@ \uline{a regular closed immersion}, @@M27@@ \uline{the conormal sheaf of} @@M28@@ \uline{in} @@M29@@. \uline{Then}, \uline{for every} @@M30@@, \uline{we have the relation}:
\sgaStableTarget{sga6.e07.eq.2-7-1}@@M31@@
\uline{where} @@M32@@ \uline{is the class of} @@M33@@ \uline{in} @@M34@@, \uline{and} @@M35@@.
